\chapter{Digital Twins for the Software Engineer}
\label{ch:dt-intro}

This chapter introduces the digital twin (DT) to a reader whose training is in
software engineering. It assumes familiarity with software architecture,
middleware, distributed systems and continuous delivery, and it assumes no
prior reading in the digital twin literature. It therefore defines the
vocabulary of that literature and leaves the vocabulary of software
engineering undefined. What the chapter does not attempt is a survey of
application domains or a comparison of commercial platforms; both are
deliberately out of scope.

The chapter argues toward one question, which the remainder of the thesis
inherits: \emph{what does treating a digital twin as a software-intensive
system, rather than as a modelling artefact, demand of software engineering
practice?} The question is worth asking because the two framings lead to
different work. Read as a modelling artefact, a twin is finished when it
predicts well. Read as a software system, it is finished when it can be
deployed, observed, changed and retired --- and prediction is one property
among several.

\section{What the term denotes}
\label{sec:dt-intro:denotation}

The concept originates in product lifecycle management, where
\citet{grieves_digital_2017} formulated it in three parts: a physical product
in real space, a virtual product in virtual space, and the data connections
tying the two together. That formulation is a useful anchor and a poor
specification, because it says nothing about what the virtual product is
\emph{for}. \citet{tao_five-dimension_2019} extend it to five dimensions by
adding data and services as first-class elements, on the argument that the
original three cannot accommodate what applications actually require. The
addition of a service layer is the step that matters here: it is the point at
which the twin acquires consumers, interfaces and availability requirements,
and therefore the point at which it becomes software.

The most operationally useful classification for a software engineer is that of
Kritzinger et al., which discriminates not by domain but by degree of data
integration between the physical and digital counterparts
\citep{kritzinger_digital_2018}. Where both directions of data flow are
manual, the artefact is a \emph{digital model}. Where the flow from physical
to digital is automatic but the return path is not, it is a \emph{digital
shadow}. Only where both directions are automatic does the term
\emph{digital twin} apply. This is a checkable property of a deployed system
rather than a claim about its ambition, which is precisely why it is worth
adopting.

It would overstate the field's maturity to present that classification as
settled. Definitional plurality is documented rather than merely suspected:
\citet{semeraro_digital_2021} survey the competing definitions without
converging on one, \citet{pfeiffer_towards_2025} are still proposing a
unifying reference model, and \citet{ellwein_rethinking_2025} report that even
within a standardised vocabulary a systematic mapping reached no consensus on
type definitions. The practical consequence for a reader of this literature is
to treat the label sceptically and read for the integration level instead. In
the manufacturing review that proposed the classification, work at the fully
integrated level was itself found to be scarce relative to work on digital
models and digital shadows \citep{kritzinger_digital_2018}.

\section{Why this is a software engineering problem}
\label{sec:dt-intro:se-problem}

The artefact an engineer must actually build, deploy and keep running is
software: models, ingestion pipelines, middleware, storage and the services
built on top. That this constitutes a software engineering concern is not this
chapter's proposal but an established framing, mapped across domains by
\citet{dalibor_cross-domain_2022} and set out as a community agenda by
\citet{cleophas_community-sourced_2022}.

Read that way, much of the construction problem is recognisable.
Architecture-level technique applies directly, and has been applied: both
\citet{tekinerdogan_systems_2020} and \citet{ferko_architecting_2022}
address twin construction through architectural patterns and views. The
integration layer has been surveyed as middleware in its own right
\citep{almeida_middleware_2023}, an open-source tooling ecosystem exists and
has been compared through case studies \citep{gil_survey_2024}, and reuse
across twins has been pursued as composition on a shared platform
\citep{talasila_composable_2025}. What the twin exists to deliver is likewise
a service catalogue --- monitoring, prediction, decision support
\citep{frasheri_advanced_2024} --- and twins are rarely deployed alone, which
makes them a systems-of-systems integration problem
\citep{michael_integration_2022, combemale_challenges_2025}.

None of this is exotic to the intended reader. That is the point of the
section: the claim is not that digital twins require a new software
engineering, but that they exercise the existing one, and that the interesting
question is where they exercise it past its usual limits.

\section{Where the familiar toolkit runs out}
\label{sec:dt-intro:limits}

Three pressure points recur, and they are where this thesis locates its
contribution.

First, correctness is defined against a physical referent that the engineer
does not control. A conventional system is correct with respect to a
specification; a twin is correct with respect to an asset that ages, is
repaired and is modified. Verification, validation and uncertainty
quantification for twins are consequently identified as open foundational
gaps rather than as solved engineering
\citep{committee_on_foundational_research_gaps_and_future_directions_for_digital_twins_foundational_2024}.
Model-driven engineering supplies part of the response, with its own
limitations acknowledged \citep{michael_model-driven_2025}.

Second, time enters the specification. A software engineer normally treats
latency as a quality attribute to be budgeted. In a twin whose value depends
on tracking a live asset, the divergence between the physical system's
timeline and the twin's is a correctness condition, and one that has to be
addressed explicitly in the design rather than absorbed
\citep{frasheri_addressing_2023}.

Third, the twin evolves continuously alongside the asset, so there is no
release after which it is stable. Evolution has accordingly been argued to sit
at the centre of digital twin engineering \citep{alskaif_evolution_2025} and
has been given explicit notational support \citep{mertens_continuous_2024}.
The consequence for practice is a convergence of modelling with operations,
pursued as model-based DevOps \citep{combemale_model-based_2023} and, for
twins specifically, as TwinOps \citep{hugues_twinops_2022}. A related and less
comfortable consequence is that much twin construction remains manual and
hard to reproduce \citep{barbie_toward_2024}, which the community has begun to
answer with shared exemplars \citep{kamburjan_greenhousedt_2024}.

\section{An open disagreement}
\label{sec:dt-intro:disagreement}

The literature does not agree on whether any of this is genuinely new.
\citet{ali_modeling_2024} pose the question directly: is the digital twin an
evolution of modelling and simulation, or a revolution? The evolutionary
reading is defensible --- co-simulation, model calibration and validation all
predate the term, and a twin assembled from them inherits their methods. This
thesis sides with that evolutionary reading as far as the modelling is
concerned, and locates the discontinuity elsewhere: what changes is the
obligation to run the model as a live, connected, evolving service against an
asset in the field. That obligation is what
sections~\ref{sec:dt-intro:se-problem} and~\ref{sec:dt-intro:limits} have
been describing, and it is enough to make the engineering problem distinct
without requiring the stronger claim that the science is.
